\section{Confining large common divisors}
\label{p13:conf:section}

At small relative rank, the host count already controls
triple probabilities. At macroscopic rank, many large
prime factors must be shared between the windows.
An interpolation argument confines the offsets of these
common divisors to a fixed-degree curve. Unless the pair
has an exact rational alignment, that curve contains too
few offsets to sustain many large groups.

Fix \(H=C_1B\) larger than the diameters and use the odd
kernel \(K_H\) from Section~\ref{p13:half:section}.
For this section, a pair is called aligned if two distinct
offset pairs give the same ratio of true values.
Equivalently, two points \((i,j)\) lie on a line through
\((-x,-y)\). Reducing the ratio bounds its height by
\(O(B)\); the enlarged height-\(B^3\) convention used
elsewhere in the paper includes these alignments.


The confinement lemma, proved in Section~\ref{p14:comp:confinement},
starts from pairwise coprime groups $g_a>N^\tau$ dividing
$x+i_a$ and $y+j_a$.  Determinants of bounded-degree monomials force
their offset points onto an algebraic curve.  Nonlinear components
have $o(B)$ grid points, while a line containing many groups must be
an exact alignment.  After removing these few groups, the surviving
coordinates have small odd kernels and the half-count machinery makes
their centres rare.  The argument never enumerates the auxiliary curve
or a chosen relation subspace.

For each fixed arity, the companion also uses an anchored version:
after coordinate regularization, each active anchor value meets only
a fixed number of minimal prime supports. Confinement removes the
large supports at a cost of $o(B)$ anchor coordinates. Two remaining
small kernels make that centre rare. On a primitive family this
works at every window, giving $N^\varepsilon$ dense hosts and the
fixed-size weights in Proposition~\ref{p13:transfer:primitive}.
The arity and density proportion are fixed before $N$.
The same section completes the stronger pair profile of
Proposition~\ref{p1:arith:hosts} without using that profile as an input.

The two forms needed to partition compatible triples are as follows.
The companion also proves the primitive-triple and dense-pair versions,
Theorem~\ref{p13:conf:dense-primitive} and
Corollary~\ref{p13:conf:dense-pairs}.

\begin{proposition}[Dense triples with weak proper intersections]
\label{p13:conf:near-primitive}
Fix \(\eta>0\) and \(0<\theta<\eta/100\).
Admissible triples satisfying
\[
 \rho_{\mathrm{rel}}(T)\ge\eta B,\qquad
 d_{xy}\le\theta B\quad\text{for every pair}
\]
have \(O_{\eta,\theta,C,\varepsilon}(N^\varepsilon)\)
hosts and total homogeneous mass
\(O(N^{-2+\theta+\varepsilon})\).
\end{proposition}

Its proof is given in Section~\ref{p14:comp:confinement:near}.

\begin{proposition}[A strong nonaligned pair in a triple]
\label{p13:conf:strong-profile}
Fix \(\theta,\sigma,\varepsilon>0\).
Compatible connected triples with a nonaligned pair
of rank at least \(\theta L\), and
\(\rho_{\mathrm{rel}}(T)\le\sigma L\), have total
joint mass
\[
 O_{\theta,\sigma,C,\varepsilon}
       \bigl(N^{-\min\{3/2,\,3-\sigma\}+\varepsilon}\bigr).
\]
\end{proposition}

Its proof is given in Section~\ref{p14:comp:confinement:strong}.



The last proposition preserves the actual conditional
probabilities. Replacing them by independent weights
would lose the large-rank case. What remains is a
pointwise bound on the possible pair ranks and a separate
calculation for rationally aligned triples.
